\appendix

\chapter{Code d’entraînement (extraits)}\label{app:code}

Les extraits ci-dessous reprennent la structure du carnet d’entraînement. La construction du modèle est reconstituée à partir de la configuration enregistrée du modèle retenu (augmentation, dropout, tête de classification)~; les valeurs des phases d’entraînement sont celles décrites au chapitre~\ref{chap:modele}.

\Needspace*{12\baselineskip}
\section{Construction du modèle}

\begin{lstlisting}[language=PythonK]
IMG, N = (224, 224), 11

augment = tf.keras.Sequential([
    tf.keras.layers.RandomFlip("horizontal_and_vertical"),
    tf.keras.layers.RandomRotation(0.15),
    tf.keras.layers.RandomZoom(0.15),
    tf.keras.layers.RandomBrightness(0.1),
    tf.keras.layers.RandomContrast(0.1),
], name="augmentation")      # pas de changement de teinte

base = tf.keras.applications.MobileNetV2(
    input_shape=IMG + (3,), include_top=False, weights="imagenet")
base.trainable = False       # phase 1 : base gelee

inputs = tf.keras.Input(shape=IMG + (3,))        # pixels 0-255
x = augment(inputs)
x = tf.keras.applications.mobilenet_v2.preprocess_input(x)   # [-1,1]
x = base(x, training=False)
x = tf.keras.layers.GlobalAveragePooling2D()(x)
x = tf.keras.layers.Dropout(0.3)(x)
outputs = tf.keras.layers.Dense(N, activation="softmax")(x)
model = tf.keras.Model(inputs, outputs)
\end{lstlisting}

\Needspace*{12\baselineskip}
\section{Poids de classes et phase 1}

\begin{lstlisting}[language=PythonK]
y_train = np.concatenate([y.numpy() for _, y in train_ds])
weights = compute_class_weight("balanced",
                               classes=np.arange(N), y=y_train)
CLASS_WEIGHT = dict(enumerate(weights))

cb1 = [tf.keras.callbacks.EarlyStopping(monitor="val_loss", patience=6,
                                        restore_best_weights=True),
       tf.keras.callbacks.ReduceLROnPlateau(monitor="val_loss",
                                            factor=0.5, patience=3)]
model.compile(optimizer=tf.keras.optimizers.Adam(1e-3),
              loss="sparse_categorical_crossentropy",
              metrics=["accuracy"])
hist1 = model.fit(train_ds, validation_data=val_ds, epochs=25,
                  class_weight=CLASS_WEIGHT, callbacks=cb1)
\end{lstlisting}

\Needspace*{12\baselineskip}
\section{Phase 2~: fine-tuning}

\begin{lstlisting}[language=PythonK]
base.trainable = True
FREEZE_UNTIL = len(base.layers) - 40          # 40 dernieres couches
for i, layer in enumerate(base.layers):
    layer.trainable = (i >= FREEZE_UNTIL) and not isinstance(
        layer, tf.keras.layers.BatchNormalization)

cb2 = [tf.keras.callbacks.EarlyStopping(monitor="val_loss", patience=10,
                                        restore_best_weights=True),
       tf.keras.callbacks.ReduceLROnPlateau(monitor="val_loss", factor=0.5,
                                            patience=4, min_lr=1e-7)]
model.compile(optimizer=tf.keras.optimizers.Adam(1e-5),
              loss="sparse_categorical_crossentropy",
              metrics=["accuracy"])
hist2 = model.fit(train_ds, validation_data=val_ds, epochs=40,
                  class_weight=CLASS_WEIGHT, callbacks=cb2)
\end{lstlisting}

Remarque~: les objets \code{cb1} et \code{cb2} sont créés séparément~; réutiliser les mêmes callbacks dans les deux phases conserverait en mémoire la meilleure perte de la phase 1 (section~\ref{sec:finetune}).

\chapter{Code de l’application Android (extraits)}\label{app:kotlin}

\Needspace*{12\baselineskip}
\section{Inférence avec test-time augmentation}

\begin{lstlisting}[language=Kotlin]
/** Probabilités des 11 classes, moyennées sur l'image et son miroir. */
private fun classify(bitmap: Bitmap): FloatArray {
    val src = bitmap.copy(Bitmap.Config.ARGB_8888, false)
    val flipped = Bitmap.createBitmap(
        src, 0, 0, src.width, src.height,
        Matrix().apply { preScale(-1f, 1f) }, true)
    val a = runModel(Bitmap.createScaledBitmap(src, IMG_SIZE, IMG_SIZE, true))
    val b = runModel(Bitmap.createScaledBitmap(flipped, IMG_SIZE, IMG_SIZE, true))
    return FloatArray(a.size) { (a[it] + b[it]) / 2f }
}

private fun runModel(scaled: Bitmap): FloatArray {
    // Pixels RGB float32 entre 0 et 255 (normalisation dans le modèle)
    val input = ByteBuffer.allocateDirect(4 * IMG_SIZE * IMG_SIZE * 3)
        .order(ByteOrder.nativeOrder())
    val pixels = IntArray(IMG_SIZE * IMG_SIZE)
    scaled.getPixels(pixels, 0, IMG_SIZE, 0, 0, IMG_SIZE, IMG_SIZE)
    for (p in pixels) {
        input.putFloat(((p shr 16) and 0xFF).toFloat())  // R
        input.putFloat(((p shr 8) and 0xFF).toFloat())   // G
        input.putFloat((p and 0xFF).toFloat())           // B
    }
    input.rewind()
    val output = Array(1) { FloatArray(labels.size) }
    interpreter.run(input, output)
    return output[0]
}
\end{lstlisting}

\Needspace*{12\baselineskip}
\section{Décision~: fiable, entre deux stades ou incertain}

\begin{lstlisting}[language=Kotlin]
val top = probs.withIndex().sortedByDescending { it.value }.take(3)
val best = labels[top[0].index]
val conf = top[0].value
val uncertain = conf < CONFIDENCE_MIN            // 0.70

val second = labels[top[1].index]
val sameFruit = best.substringBefore("_") == second.substringBefore("_")
val between = uncertain && sameFruit && (conf + top[1].value) >= 0.85f

if (between) {
    // deux stades voisins du même fruit se partagent le résultat
} else if (uncertain) {
    // "Incertain" : meilleure hypothèse donnée à titre indicatif
} else {
    // résultat fiable : classe, confiance et conseil
}
\end{lstlisting}

\chapter{Résultats détaillés}\label{app:resultats}

\Needspace*{12\baselineskip}
\section{Classes et ordre des sorties}

L’ordre des 11 sorties du modèle (fichier \code{labels.txt}) est l’ordre alphabétique suivant~:

\begin{enumerate}[nosep,leftmargin=*]
  \item \code{avocat\_mur}
  \item \code{avocat\_non\_mur}
  \item \code{avocat\_partiellement\_mur}
  \item \code{avocat\_pourri}
  \item \code{banane\_mure}
  \item \code{banane\_non\_mure}
  \item \code{banane\_trop\_mure}
  \item \code{fraise\_mure}
  \item \code{fraise\_non\_mure}
  \item \code{fraise\_partiellement\_mure}
  \item \code{fraise\_pourrie}
\end{enumerate}

\Needspace*{12\baselineskip}
\section{Rappel des chiffres clés}

\begin{table}[H]
\centering
\caption{Chiffres clés du projet.}\label{tab:cles}
\small
\begin{tabular}{@{}lr@{}}
\toprule
\textbf{Grandeur} & \textbf{Valeur} \\
\midrule
Photographies d’origine & 1\,431 \\
Images de fruits recadrées & 1\,598 \\
Entraînement / validation / test & 1\,118 / 249 / 231 \\
Paramètres du modèle & 2\,272\,075 \\
Exactitude de test (Keras) & 93,07\,\% (215/231) \\
Intervalle de Wilson à 95\,\% & [89,0\,\% ; 95,7\,\%] \\
F1 macro & 0,897 \\
Exactitude de test (TFLite quantifié) & 92,64\,\% (214/231) \\
Taille Keras / TFLite & 23,1\,Mo / 2,5\,Mo \\
Inférence CPU (Kaggle) & 19,6\,ms \\
\bottomrule
\end{tabular}
\end{table}
